\chapter{Rückkopplung durch Photonen und Rotverschiebung}

Die Rotverschiebung ist in der WDBT+ kein Doppler-Effekt, sondern eine \emph{kumulative gravitative Energieabgabe}:

\begin{equation}
z(d) = \frac{1}{c^2} \int_0^d \frac{G M(r)}{r^2} dr.
\end{equation}

Die zurückfließende Energie ist die Energie der Rotverschiebung. Jedes Photon – egal ob von Supernova, Stern oder Planetenkern – gibt auf seinem Weg durch das Universum Energie an das Vakuum ab.

\subsection{Polyrhythmisches System}

Es gibt unzählige Taktgeber:

\begin{itemize}
    \item Supernovae (hochenergetische Photonen),
    \item Sterne (optische Photonen),
    \item Planetenkerne (Infrarot-Photonen).
\end{itemize}

Der Kosmos ist ein \emph{polyrhythmisches System} – ein Netzwerk von Oszillatoren, die alle zur Rückkopplung beitragen.

\subsection{Die Rolle des Quantenpotentials Q}

Die Rückkopplung wird durch das Quantenpotential Q reguliert. Q ist instantan und nicht-lokal – es kann Energie von einem Ort zum anderen transferieren, ohne dass ein Signal mit \(c\) unterwegs sein muss.

Die Energieabgabe der Photonen durch Rotverschiebung wird von Q aufgenommen und in neue Materie umgewandelt (\( \sigma_{\text{cre}} \)).